% =============================================================================
% definitions.tex — formal definitions of the streaming framework
%
% Single source of truth for the math behind backend-modules/streaming.
% Every work package of issue #1 that introduces or changes a concept updates
% this file in the same commit. Code docstrings refer to the definitions here
% by label (e.g. "see def:run"). Sections marked TODO are owned by the work
% package named there.
%
% Requires: amsmath, amssymb, amsthm (definition, proposition, remark envs).
% =============================================================================

% --- universes (as in the preliminaries) -------------------------------------
\providecommand{\Usigma}{\mathbb{U}_{\Sigma}}
\providecommand{\Uact}{\mathbb{U}_{\mathit{act}}}
\providecommand{\Uobj}{\mathbb{U}_{\mathit{obj}}}
\providecommand{\Utype}{\mathbb{U}_{\mathit{type}}}
\providecommand{\Uattr}{\mathbb{U}_{\mathit{attr}}}
\providecommand{\Uval}{\mathbb{U}_{\mathit{val}}}
\providecommand{\Utime}{\mathbb{U}_{\mathit{time}}}
\providecommand{\Uqual}{\mathbb{U}_{\mathit{qual}}}

% --- framework notation ------------------------------------------------------
\newcommand{\dom}{\mathit{dom}}
\newcommand{\upd}{\mathit{upd}}
\newcommand{\fin}{\mathit{fin}}
\newcommand{\touched}{\mathit{touched}}
\newcommand{\score}{\mathit{score}}
\newcommand{\dl}{\delta}                    % deadline function
\newcommand{\hupd}{\rho}                    % heuristic state update
\newcommand{\proj}[2]{#1\!\Downarrow_{#2}}  % type projection O⇂t
\newcommand{\obj}{\mathit{obj}}
\newcommand{\otype}{\mathit{type}}
\newcommand{\alast}{a_{\mathit{last}}}
\newcommand{\Rcap}{\textsc{cap}}
\newcommand{\Rexp}{\textsc{exp}}
\newcommand{\Rdis}{\textsc{dis}}

% =============================================================================
\section{Object-Centric Event Streams}\label{sec:streams}
% =============================================================================

Small adaptation of the streaming definition of~\cite{Liss2025}.

\begin{definition}[Object-centric event stream]\label{def:stream}
An \emph{object-centric event stream} is a possibly infinite sequence
$S = \langle s_1, s_2, \dots \rangle$ of items $s_i = (x_i, t_i)$ with a
timestamp $t_i \in \Utime$, where $x_i$ is an event item, an object-to-object
update or an attribute update.
\begin{itemize}
  \item An \emph{event item} is a tuple $x = (a, R, A)$ with the activity
        $a \in \Uact$, a finite set of typed, optionally qualified E2O links
        $R \subseteq \Uobj \times \Utype \times (\Uqual \cup \{\bot\})$, and the
        event attributes $A \colon \Uattr \rightharpoonup \Uval$.
        Its object set is $\obj(x) = \{o \mid (o, \cdot, \cdot) \in R\}$.
  \item An \emph{object-to-object update} is a tuple
        $x = (o_1, ot_1, q, o_2, ot_2)$ with $o_1, o_2 \in \Uobj$,
        $ot_1, ot_2 \in \Utype$ and $q \in \Uqual \cup \{\bot\}$.
  \item An \emph{attribute update} is a tuple
        $x = (o, ot, \mathit{att}, v) \in \Uobj \times \Utype \times \Uattr \times (\Uval \cup \{\bot\})$,
        where $v = \bot$ clears the attribute.
\end{itemize}
We write $\mathcal{I}$ for the set of all stream items and
$\mathcal{I}_{\mathit{ev}} \subseteq \mathcal{I}$ for the event items.
The stream is \emph{well-formed} if
(i)~$t_i \le t_j$ for all $i \le j$, and
(ii)~if $o$ occurs with type $ot$ in $s_i$ and with type $ot'$ in $s_j$, then
$ot = ot'$; we write $\otype(o)$ for this type.
\end{definition}

\paragraph{Notation.}
The index $i$ of $s_i$ is its \emph{position}; all ages, deadlines and
thresholds in this document are measured in positions.
For a set of objects $O$ and a type $t$, $\proj{O}{t} = \{o \in O \mid \otype(o) = t\}$.
$S_{\le i} = \langle s_1, \dots, s_i \rangle$ is the prefix up to position $i$.

\begin{remark}[Positions vs.\ events]
If the stream contains only event items (the setting of the online OCPT
miner), positions coincide with event indices, so $\Delta(o)$ of
Def.~\ref{def:lifecycle} is exactly the inactivity of~\cite{Muehlmeyer2026}.
\end{remark}

\begin{remark}[Replay convention, evaluation only]
When an OCEL is replayed as a stream, items are ordered by timestamp, ties by
kind (events, attribute updates, O2O updates) and then by identifier. O2O
relations carry no timestamp in OCEL~2.0; each is placed at the first event of
its source object, as in~\cite{Liss2025}.
\end{remark}

% =============================================================================
\section{Buffer-Based Streaming Framework}\label{sec:framework}
% =============================================================================

A streaming algorithm cannot keep everything it has seen. Everything an
algorithm remembers is stored in buffers. Each buffer has an optional capacity
and an eviction heuristic. Algorithms differ in what they store; eviction
strategies differ only in their heuristic.

\begin{definition}[Buffer]\label{def:buffer}
A buffer $B = (K, V, c, H)$ stores values from $V \subseteq \Usigma$ under keys
from $K \subseteq \Usigma$ (e.g.\ objects, arcs, activities). Its content is a
partial map $b \colon K \rightharpoonup V$. The capacity is
$c \in \mathbb{N} \cup \{\infty\}$ and $H$ is an eviction heuristic
(Def.~\ref{def:heuristic}). For each key $k \in \dom(b)$ the buffer records
its \emph{bookkeeping}: the position $f(k)$ of its first write, the position
$l(k)$ of its last write, and its number of writes $n(k)$.
\end{definition}

\begin{definition}[Eviction heuristic]\label{def:heuristic}
An \emph{eviction heuristic} $H = (\Theta, \theta_0, \score, \dl, \hupd)$
consists of
\begin{itemize}
  \item a set of heuristic states $\Theta$ with initial state $\theta_0$;
  \item a \emph{score} $\score(k, \sigma, \theta, i) \in \mathbb{R}$; a higher
        score means ``evict sooner'';
  \item a \emph{deadline} $\dl(k, \sigma, \theta, i) \in \mathbb{N} \cup \{\infty\}$,
        the position at which $k$ should be evicted;
  \item a \emph{state update} $\hupd(\theta, k, \mathit{op}, \sigma, i) \in \Theta$
        applied whenever its buffer writes or evicts $k$, with
        $\mathit{op} \in \{\mathit{write}\} \cup \{\mathit{evict}(r) \mid r \in \{\Rcap, \Rexp, \Rdis\}\}$.
\end{itemize}
A heuristic need not define both $\score$ and $\dl$; it supports
\emph{capacity mode} if it defines $\score$ and \emph{expiry mode} if it
defines $\dl$. Stateless heuristics have $|\Theta| = 1$.
\end{definition}

The heuristic state $\theta$ belongs to the heuristic, not to the algorithm, so
any heuristic can be combined with any algorithm. It counts toward memory
(Rem.~\ref{rem:memory}).

\begin{definition}[Streaming algorithm]\label{def:algorithm}
A \emph{streaming algorithm} $A = (\mathcal{B}, \upd, \fin, \lambda)$ consists
of a tuple of buffers $\mathcal{B} = (B_1, \dots, B_n)$ and three functions on
algorithm states $\sigma = (b_1, \dots, b_n)$, initially all empty:
\begin{itemize}
  \item $\upd(\sigma, s)$ processes item $s$ and returns the new state;
        $\touched_j(s) \subseteq K_j$ are the keys of $B_j$ it \emph{reads or
        writes};
  \item $\fin(\sigma, j, k)$ is called after key $k$ was removed from $B_j$ and
        returns the new state; it may change other buffers (e.g.\ record the
        end activity of an evicted object, or discard dependent keys);
        keys it writes are added to $\touched$;
  \item $\lambda(\sigma)$ is the current model.
\end{itemize}
We require that the ``$\fin$ of $B_j$ writes to $B_l$'' relation is acyclic.
\end{definition}

\begin{definition}[Run]\label{def:run}
Each buffer $B_j$ additionally keeps a schedule
$d_j \colon \dom(b_j) \to \mathbb{N} \cup \{\infty\}$ (only if $H_j$ is in
expiry mode). For each item $s_i$:
\begin{enumerate}
  \item \textbf{Update.} $\sigma \leftarrow \upd(\sigma, s_i)$. For every key
        $k$ written in $B_j$: update its bookkeeping,
        $\theta_j \leftarrow \hupd_j(\theta_j, k, \mathit{write}, \sigma, i)$,
        and, in expiry mode, $d_j(k) \leftarrow \dl_j(k, \sigma, \theta_j, i)$.
  \item \textbf{Evict} until no rule applies, buffers in order $j = 1, \dots, n$
        and repeated as long as some $\fin$ wrote to an earlier buffer.
        Let $C_j = \dom(b_j) \setminus \touched_j(s_i)$ be the candidates.
        \begin{enumerate}
          \item \emph{Expiry} (if $H_j$ is in expiry mode): while some
                $k \in C_j$ has $d_j(k) \le i$, take the one with the smallest
                $d_j(k)$ (ties: smallest $f(k)$) and set
                $d' = \dl_j(k, \sigma, \theta_j, i)$.
                If $d' > i$, reschedule $d_j(k) \leftarrow d'$;
                otherwise evict $k$ with reason $\Rexp$.
          \item \emph{Capacity} (if $c_j < \infty$): while
                $|\dom(b_j)| > c_j$ and $C_j \neq \emptyset$, evict
                $k^* = \arg\max_{k \in C_j} \score_j(k, \sigma, \theta_j, i)$
                (ties: smallest $f(k)$) with reason $\Rcap$.
        \end{enumerate}
        \emph{Evicting} $k$ from $B_j$ with reason $r$ means: remove $k$ from
        $b_j$ and $d_j$, then
        $\sigma \leftarrow \fin(\sigma, j, k)$ and
        $\theta_j \leftarrow \hupd_j(\theta_j, k, \mathit{evict}(r), \sigma, i)$.
        Keys written by $\fin$ are handled as in step~1 (bookkeeping, $\hupd$
        with $\mathit{write}$, schedule) before $\hupd_j$ is applied.
        An algorithm may also remove a key itself (reason $\Rdis$, e.g.\ from
        $\fin$); this calls $\hupd_j$ but not $\fin$.
\end{enumerate}
The model after item $s_i$ is $\lambda(\sigma_i)$. The \emph{unbounded run}
uses $c_j = \infty$ and no expiry for all buffers; it never evicts.
\end{definition}

The schedule realizes ``re-check when the deadline is reached'': a key is
evicted no earlier than its scheduled deadline, and only if the deadline
recomputed at that moment is not in the future. If a deadline \emph{shrinks}
after it was scheduled, the key still waits until the old deadline. This is
the semantics of~\cite{Muehlmeyer2026} and its reference implementation.

Termination: step~2 terminates because every eviction removes a key that is
not re-added in the same step (written keys are touched), and the $\fin$
relation is acyclic.

\begin{definition}[Premature eviction]\label{def:premature}
A key $k$ evicted from $B_j$ at item $s_i$ is evicted \emph{too early} if
$k \in \touched_j(s_{i'})$ for some $i' > i$.
\end{definition}

\begin{proposition}[Exactness]\label{prop:exact}
Assume
(i)~evicting a key commutes with processing any item that does not touch it, and
(ii)~evicting a key does not change the current model,
$\lambda(\fin(\sigma, j, k)) = \lambda(\sigma)$.
If no key is evicted too early, the run yields the same model as the
unbounded run.
\end{proposition}
\begin{proof}
No later item touches an evicted key. By (i), every eviction can be moved
step by step to the end of the run without changing the final state. By (ii),
the evictions at the end do not change the model. What remains is the
unbounded run.
\end{proof}

\begin{remark}[Where (i) and (ii) hold]\label{rem:exactness}
(i) relies on $\touched$ containing read keys; all building blocks of
Sec.~\ref{sec:blocks} only read keys they also write.
(ii) holds for object buffers whose $\fin$ moves the object's contribution into
the model (Def.~\ref{def:df}: its end activity). It fails for buffers that
store the model itself: Löseke's arc observation log, and equally the counter
buffers of Def.~\ref{def:df} when bounded. Bounding such buffers always costs
accuracy; bounding object buffers costs accuracy only when the heuristic
evicts too early.
\end{remark}

\begin{remark}[Memory]\label{rem:memory}
After each item, $|\dom(b_j)| \le \max(c_j, |\touched_j(s_i)|)$ for every
buffer; $c_j \ge \max_i |\touched_j(s_i)|$ (e.g.\ the maximal number of objects
per event for an object buffer) gives $|\dom(b_j)| \le c_j$.
Memory is bounded if every buffer whose size grows with the number of objects
has a finite capacity and every heuristic state $\theta_j$ is bounded
independently of the number of objects. Memory is measured in buffer entries.
\end{remark}

% =============================================================================
\section{Building Blocks}\label{sec:blocks}
% =============================================================================

Blocks are reusable parts of algorithms. Each owns buffers and defines its
share of $\upd$, $\fin$ and $\lambda$. Unless stated otherwise, consider an
event item $s_i = ((a, R, A), t_i)$ with $O = \obj(x_i)$.

\begin{definition}[Object tracker]\label{def:tracker}
Buffer $B_{\mathit{obj}}$ with keys $\Uobj$ and values $\Uact \times \Utime$
(last activity and its timestamp). On $s_i$, for each $o \in O$ (in relation
order): emit $\mathit{transition}(o, b, t', a, t_i, |\proj{O}{\otype(o)}|)$ where
$(b, t') = b_{\mathit{obj}}(o)$, or $(\bot, \bot)$ if $o \notin \dom(b_{\mathit{obj}})$,
then write $b_{\mathit{obj}}(o) \leftarrow (a, t_i)$. We write $\alast(o)$ for
the activity in $b_{\mathit{obj}}(o)$; $l(o)$ is the position of $o$'s last
event. Evicting $o$ emits $\mathit{retired}(o, \alast(o))$.
\end{definition}

\begin{definition}[Directly-follows counters]\label{def:df}
Buffers $B_{\mathit{arcs}}$ with keys $\Uact \times \Uact \times \Utype$,
$B_{\mathit{start}}$ and $B_{\mathit{end}}$ with keys $\Uact \times \Utype$,
values $\mathbb{N}$.
$\mathit{transition}(o, \bot, \cdot, a, \cdot, \cdot)$ increments $b_{\mathit{start}}(a, \otype(o))$;
$\mathit{transition}(o, b, \cdot, a, \cdot, \cdot)$ with $b \ne \bot$ increments
$b_{\mathit{arcs}}(b, a, \otype(o))$;
$\mathit{retired}(o, a)$ increments $b_{\mathit{end}}(a, \otype(o))$.
The end activities of type $t$ are taken from \emph{all} objects:
\[
  \mathit{End}(a, t) = b_{\mathit{end}}(a, t) + |\{o \in \dom(b_{\mathit{obj}}) \mid \otype(o) = t, \alast(o) = a\}|.
\]
\end{definition}
This is what makes Rem.~\ref{rem:exactness}(ii) hold for $B_{\mathit{obj}}$.
(The Rust reference of~\cite{Muehlmeyer2026} uses only the second term.)

\begin{definition}[Activity--type relations]\label{def:relations}
Buffers $B_{\mathit{act}}$ (keys $\Uact$) and $B_{\mathit{rel}}$ (keys
$\Uact \times \Utype$) hold the counters $N(a)$ (events of $a$), $N(a,t)$ (events of $a$ with
$\proj{O}{t} \ne \emptyset$) and $N_{>1}(a,t)$ (events of $a$ with
$|\proj{O}{t}| > 1$). Then
$\mathit{rel} = \{(a,t) \mid N(a,t) > 0\}$,
$\mathit{con} = \{(a,t) \mid N_{>1}(a,t) > 0\}$,
$\mathit{def} = \{(a,t) \in \mathit{rel} \mid N(a,t) < N(a)\}$.
\end{definition}

\begin{definition}[Divergence index]\label{def:divergence}
Let $h \colon 2^{\Uobj} \to \mathbb{Z}$ be a fast hash
(implementation: Python \texttt{hash(frozenset(ids))}; collisions accepted).
Buffer $B_{\mathit{div}}$ with keys $\Uact \times \Utype \times \mathbb{Z}$ and
values $\mathbb{Z}$, and a buffer $B_{\mathit{divergent}}$ holding the set
$\mathit{div} \subseteq \Uact \times \Utype$ (small; normally unbounded). On $s_i$, for each type $t$ with
$\proj{O}{t} \ne \emptyset$ and $(a,t) \notin \mathit{div}$, let
$k = (a, t, h(\proj{O}{t}))$:
if $k \in \dom(b_{\mathit{div}})$ and $b_{\mathit{div}}(k) \ne h(O)$, add
$(a,t)$ to $\mathit{div}$ and discard all keys $(a, t, \cdot)$ (reason $\Rdis$);
otherwise write $b_{\mathit{div}}(k) \leftarrow h(O)$.
We write $\mathit{div}(a) = \{t \mid (a,t) \in \mathit{div}\}$.
\end{definition}

\begin{definition}[DF\textsuperscript{2} graph]\label{def:df2}
From Defs.~\ref{def:df}--\ref{def:divergence}, the DF\textsuperscript{2} graph
is $(\Sigma, E_\to, \Sigma_{\mathit{start}}, \Sigma_{\mathit{end}})$ with
$\Sigma = \{a \mid N(a) > 0\}$ and
\begin{align*}
  E_\to &= \{(a,b) \mid \exists t \colon b_{\mathit{arcs}}(a,b,t) > 0 \wedge t \notin \mathit{div}(a) \cap \mathit{div}(b)\},\\
  \Sigma_{\mathit{start}} &= \{a \mid \exists t \colon b_{\mathit{start}}(a,t) > 0 \wedge t \notin \mathit{div}(a)\},\\
  \Sigma_{\mathit{end}} &= \{a \mid \exists t \colon \mathit{End}(a,t) > 0 \wedge t \notin \mathit{div}(a)\}.
\end{align*}
The divergence filter on start and end activities is a design decision; the
offline reference applies the same filter.
The offline reference computes the same graph from the complete log, with
events in replay order, exact object sets instead of hashes, and
$\mathit{End}(a,t)$ from the last activity of every object.
\end{definition}

\begin{definition}[OCPT discovery]\label{def:ocpt}
A \emph{cut miner} maps a DF\textsuperscript{2} graph to an object-centric
process tree. The default cut miner tries the cuts xor, sequence,
concurrency and loop in this order, recursing on the parts. If no cut applies
to an activity set $X$ with $|X| > 1$, it returns the node
$\mathit{NoCut}(X)$; there is no fallback. Each leaf $a$ is annotated with
$(a, \mathit{rel}(a), \mathit{con}(a), \mathit{div}(a), \mathit{def}(a))$
from Defs.~\ref{def:relations} and~\ref{def:divergence}.
\end{definition}

\paragraph{Cut conditions (default cut miner).}
Let $X$ be the current activity set, $E_X = E_\to \cap (X \times X)$, and
$S_X, F_X \subseteq X$ its start and end activities: activities entered from
(left to) an activity outside $X$ in the whole graph, plus the start (end)
activities of the parent. All cuts are binary, $X = X_1 \uplus X_2$.
\begin{itemize}
  \item \emph{xor}: $X_1$ is the connected component (undirected $E_X$) with
        the smallest activity, $X_2$ the rest; no edge connects them.
  \item \emph{sequence}: from the condensation of $E_X$ into strongly
        connected components, $X_1$ holds the components with outgoing and
        $X_2$ those with incoming condensation edges (components with both are
        assigned to $X_2$ if every $X_1$ component reaches but is not reached
        by them, else to $X_1$); every $a \in X_1$ reaches every $b \in X_2$ in
        $E_X$ and no $b$ reaches an $a$.
  \item \emph{parallel}: greedily in activity order, every $a \in X_1$ and
        $b \in X_2$ have $(a,b), (b,a) \in E_X$; both parts contain a start and
        an end activity.
  \item \emph{loop}: $X_1 \supseteq S_X \cup F_X$ (body); every other
        activity is reachable without passing an end activity from $S_X$
        (body) or without passing a start activity from $F_X$ (redo), but not
        both; every end activity has an edge into $X_2$ and every start
        activity an edge from $X_2$.
\end{itemize}
This is a port of \texttt{find\_cuts\_start} of the Rust reference and yields
the same trees wherever a cut exists (tested on 801 graphs).

% TODO(WP7): Liss et al. building blocks — arc observation log
%   (key: fresh id; value (b|⊥, a, t, o, 1/|O⇂t|, duration)), relation log,
%   object interactions (log and event cardinalities), TOTeM construction,
%   coupled removal.

% =============================================================================
\section{Eviction Heuristics}\label{sec:heuristics}
% =============================================================================

Stateless heuristics, using the bookkeeping of Def.~\ref{def:buffer}
(capacity mode):
\begin{align*}
  \text{FIFO: } & \score(k) = i - f(k) &
  \text{LIFO: } & \score(k) = f(k) \\
  \text{LRU: } & \score(k) = i - l(k) &
  \text{LFU: } & \score(k) = -n(k) \\
  \text{Random: } & \score(k) \sim U(0,1) \text{ (seeded)}
\end{align*}
Max age (expiry mode), with age limit $T$: $\dl(k) = l(k) + T$.

\begin{definition}[Lifecycle heuristic~\cite{Muehlmeyer2026}]\label{def:lifecycle}
For an object buffer (Def.~\ref{def:tracker}), with parameters
$0 \le T_{\mathit{hint}} < T_{\mathit{max}}$,
$z = 1$ and $\beta = 0.95$.
The state is an end-activity histogram $\theta = (\kappa, \nu)$ with
$\kappa \colon \Uact \times \Utype \to \mathbb{N}$ and
$\nu \colon \Utype \to \mathbb{N}$, initially zero.
For $k = \kappa(a,t)$ and $n = \nu(t)$, the confidence that $a$ ends the
lifecycle of an object of type $t$ is
\[
  C(a,t) = \begin{cases}
    0 & n = 0\\
    \max(0, \mu - z s) & \text{otherwise}
  \end{cases}
  \quad
  \mu = \frac{k+1}{n+2},\;
  s = \sqrt{\frac{(k+1)(n-k+1)}{(n+2)^2(n+3)}},
\]
and the \emph{eviction age}
\[
  \alpha(a,t) = \begin{cases}
    T_{\mathit{max}} & C(a,t) \le 0\\
    T_{\mathit{hint}} + 1 & C(a,t) \ge 0.9999\\
    T_{\mathit{hint}} + \bigl\lceil (T_{\mathit{max}} - T_{\mathit{hint}}) \cdot \beta^{C/(1-C)} \bigr\rceil & \text{otherwise.}
  \end{cases}
\]
With $\Delta(o) = i - l(o)$:
\begin{itemize}
  \item deadline: $\dl(o) = l(o) + \alpha(\alast(o), \otype(o))$;
  \item state update on $\mathit{evict}(r)$ of $o$ at position $i$: if
        $r = \Rcap$, or $r = \Rexp$ and $\Delta(o) \ge T_{\mathit{max}}$, then
        $\kappa(\alast(o), \otype(o)) \mathrel{+}= 1$ and
        $\nu(\otype(o)) \mathrel{+}= 1$;
  \item no score: in this version the heuristic supports expiry mode only.
\end{itemize}
\end{definition}

This is the heuristic of~\cite{Muehlmeyer2026}, checked after every item.
(The Rust reference batches its checks every $10\,000$ events; we do not.)
The divergence index of the online OCPT miner uses max age with
$T = T_{\mathit{max}}$.
Counting capacity evictions in the histogram has no effect while the object
buffer is unbounded; it fixes the semantics for later heuristics that combine
the lifecycle schedule with a fixed capacity.

% TODO(WP6): LFU-DA, the eight priority features of Liss et al.,
%   min-max normalized rank combination (weighted heuristic).
